% Tutorial set: 05
% Source: tut5.pdf, Tutorial Five (Process Synchronization - Part II), pp. 5-1--5-2 (PDF pp. 1--2), Questions 1--3
% Reference solution: not provided; answers checked against the standard semaphore semantics from Lectures 9--10
% Session date: 2026-09-15
% Human verified: Yes

\clearpage
\subsection{Tutorial 05：Process Synchronization Part II}
\label{tutorial:SC2005-TUT05}

\exercisemeta
  {SC2005-TUT05}
  {2026-09-15}
  {tut5.pdf，pp.~5-1--5-2（PDF pp.~1--2），Questions 1--3；原文件未附参考答案}
  {三题均已作答，并依据 Lecture 9--10 的 semaphore 标准语义完成订正}
  {已确认（2026-09-15）}

\exercisequestion{SC2005-TUT05-Q01}{两个 Basket Semaphores 的 Deadlock-free Acquisition}

\textbf{对应讲义：}
\relatedtopic{SC2005-L10-T01}{Semaphore 的系统化设计}；
\relatedtopic{SC2005-L10-T02}{Deadlock 与 Starvation}

\subsubsection*{题目}

Apples 与 oranges 在随机时刻分别进入两个 baskets。Alice 每次需要从对应 baskets 取出恰好 1 个 apple 与 2 个 oranges 才能配对得分。Apple basket 与 orange basket 分别由 semaphore $A$ 与 $O$ 保护；访问 basket 前必须取得对应 semaphore。

只允许在下列五个框中放入任意数量（可为零）的 \texttt{Wait(A)}、\texttt{Wait(O)}、\texttt{Signal(A)}、\texttt{Signal(O)}。补全程序，并保证即使其他 players 采用不同代码，Alice 的 acquisition pattern（获取模式）也不会参与形成 deadlock。

\begin{verbatim}
[1]
// Pick 1 apple

[2]
if (at least 2 oranges in basket) {
    // Pick 2 oranges and pair with apple
[3]
} else {
[4]
    // Drop 1 apple back in basket
[5]
}
\end{verbatim}

\begin{checkedsolution}
一种不同时持有两个 basket locks 的填写为
\begin{verbatim}
[1] Wait(A)

// Pick 1 apple

[2] Signal(A)
    Wait(O)

if (at least 2 oranges in basket) {
    // Pick 2 oranges and pair with apple
[3] Signal(O)
} else {
[4] Signal(O)
    Wait(A)

    // Drop 1 apple back in basket
[5] Signal(A)
}
\end{verbatim}

Alice 取出 apple 后可暂时私有地持有它，不必继续锁住 apple basket。她先释放 $A$，之后才等待 $O$；若 oranges 不足，又先释放 $O$，再等待 $A$ 并放回 apple。因此 Alice 在任何 \texttt{Wait} 处都不持有另一 basket semaphore，不可能成为 circular wait（循环等待）中的 ``hold one, wait for another'' 节点。

该结论假定其他 players 仍遵守 semaphore protocol（最终释放已取得的锁）。若其他代码永久不释放 semaphore，或多个其他 players 已经彼此 deadlock，Alice 可能被间接阻塞；任何只修改 Alice 本地代码的方案都无法排除这种外部故障。
\end{checkedsolution}

\exercisequestion{SC2005-TUT05-Q02}{Shifted Semaphore Value 与 Off-by-one Conditions}

\textbf{对应讲义：}
\relatedtopic{SC2005-L09-T02}{Binary Semaphore 实现 Mutual Exclusion}；
\relatedtopic{SC2005-L09-T04}{Blocking Semaphore：Value 与 Waiting Queue}

\subsubsection*{题目}

假设下列 \texttt{Wait} 与 \texttt{Signal} system calls（系统调用）均为 atomic，并以初值 $S.value=0$ 的 semaphore 保护每个 process 的 critical section。判断该实现是否保证 mutual exclusion（互斥）与 progress（推进性），并说明理由。

\begin{verbatim}
Wait(S):                         Signal(S):
    S.value--                        S.value++
    if (S.value < -1)                if (S.value <= 0)
        block()                          wakeup()

Process Pi:
    while (1) {
        Wait(S)
        critical section
        Signal(S)
        remainder section
    }
\end{verbatim}

\begin{checkedsolution}
该实现把常见初值为 $1$ 的 binary semaphore value 整体下移一位。其状态含义为
\[
  S.value=0: \text{无 holder、无 waiter},
\]
\[
  S.value=-(k+1):\ \text{一名 holder 与 $k$ 名 waiting processes}.
\]
前三个 processes 依次执行 \texttt{Wait} 时：
\begin{center}
\small
\begin{tabular}{@{}lcl@{}}
  \toprule
  Operation & New value & Result \\
  \midrule
  $P_1:\ \texttt{Wait(S)}$ & $-1$ & 不阻塞，进入 critical section \\
  $P_2:\ \texttt{Wait(S)}$ & $-2$ & $-2<-1$，进入 waiting queue \\
  $P_3:\ \texttt{Wait(S)}$ & $-3$ & $-3<-1$，进入 waiting queue \\
  \bottomrule
\end{tabular}
\end{center}
因此任意时刻只有第一个未阻塞的 process 位于 critical section，\textbf{mutual exclusion 成立}。

若 $P_1$ 离开时有两个 waiters，\texttt{Signal} 使
\[
  -3\longrightarrow-2,
\]
条件 $S.value\le0$ 成立，系统唤醒一个 waiter；之后每次 holder 释放许可都会继续唤醒下一名 waiter。因此在 scheduler 与 waiting queue 正常工作的前提下，\textbf{progress 成立}。

边界问题发生在没有 waiters 时：
\[
  -1\xrightarrow{\texttt{Signal}}0,
\]
代码仍会执行 \texttt{wakeup()}，但 queue 为空。若 conventional contract 规定 ``对空队列调用 \texttt{wakeup} 是安全的 no-op''，这只是多余调用，不破坏两个性质；若该行为未定义，则题目不足以给出无条件安全保证。对于本题的 shifted representation，更精确的唤醒条件是
\[
  \boxed{S.value<0}.
\]
\end{checkedsolution}

\exercisequestion{SC2005-TUT05-Q03}{用 TestAndSet 保护 Blocking Semaphore Internals}

\textbf{对应讲义：}
\relatedtopic{SC2005-L08-T04}{Synchronization Hardware 与 TestAndSet}；
\relatedtopic{SC2005-L09-T04}{Blocking Semaphore：Value 与 Waiting Queue}；
\relatedtopic{SC2005-L09-T05}{Blocking 实现的原子边界与选择}

\subsubsection*{题目}

给定
\begin{verbatim}
typedef struct {
    int value;
    struct process *L;
} semaphore;
\end{verbatim}
说明如何用 TestAndSet instruction（测试并置位指令）实现 semaphore 的 \texttt{Wait(S)} 与 \texttt{Signal(S)}。注意 $S.value$ 和 process queue $S.L$ 都是多个 processes 共享的变量。

\begin{checkedsolution}
原结构缺少保护内部状态的 lock。实现时必须为每个 semaphore 增加一个初值为 \texttt{false} 的 \texttt{guard}，或在结构外维护一一对应的 guard：
\begin{verbatim}
typedef struct {
    int value;
    struct process *L;
    bool guard;              // false: unlocked, true: locked
} semaphore;
\end{verbatim}
采用标准 TestAndSet 语义：它原子地返回旧值并把 target 写为 \texttt{true}。内部 guard 的 acquire/release 为
\begin{verbatim}
acquire_guard(S):
    while (TestAndSet(&S.guard))
        ;                   // short spin only

release_guard(S):
    S.guard = false
\end{verbatim}

Blocking semaphore 可写为
\begin{verbatim}
Wait(S):
    acquire_guard(S)
    S.value--

    if (S.value < 0) {
        enqueue(S.L, current_process)
        block_and_release(current_process, &S.guard)
    } else {
        release_guard(S)
    }

Signal(S):
    acquire_guard(S)
    S.value++

    if (S.value <= 0) {
        P = dequeue(S.L)
        wakeup(P)
    }

    release_guard(S)
\end{verbatim}

\texttt{guard} 只保护对 $S.value$ 与 $S.L$ 的短暂更新；等待 resource permit（资源许可）的 process 应进入 blocked state，而不能一直持有 guard 或长期 spinning。

\texttt{Wait} 必须先把 \texttt{current\_process} 加入 $S.L$，再阻塞；若先 \texttt{block()}，后面的 enqueue 根本不会在睡眠前执行。更细致地说，也不能简单写成
\begin{verbatim}
release_guard(S)
block()
\end{verbatim}
因为 \texttt{Signal} 可能在两句之间执行：它尝试唤醒尚未 blocked 的 process，随后该 process 才睡眠，造成 lost wakeup（唤醒丢失）。真实 kernel 需要 \texttt{block\_and\_release} 原子协议，或在 guard 保护下先 enqueue 并把 process state 标为 blocked，再释放 guard 并调用 scheduler。

同理，\texttt{Signal} 应从 $S.L$ dequeue 一个具体 process 再 \texttt{wakeup(P)}；向 queue 中加入元素属于 \texttt{Wait} 的职责。
\end{checkedsolution}

\subsubsection*{本次 Tutorial 收口}

Q1 说明跨多个 resources 的安全代码不只需要 ``最终释放''，还要检查每个 blocking point 当时持有哪些 locks；Q2 展示 semaphore value 可以采用不同数值编码，但 block/wakeup 阈值必须与状态 invariant 一致；Q3 区分用于短暂保护 semaphore internals 的 TestAndSet guard 与等待 resource permit 的 blocking queue，并揭示 enqueue、state transition、unlock 和 scheduling 之间的 lost-wakeup race。
